\bookchapter{附錄F：摺疊屏選購與使用專題}

\chapterintro{摺疊屏已經從“未來科技”變成“現在能買”的成熟品類。但它和直板機在形態、使用習慣、保養方式上差異很大，值不值得買、怎麼選、怎麼用，這一篇幫你理清楚。}

\booksection{一、摺疊屏現在發展到什麼程度了？}

先說結論：\textbf{摺疊屏已經不是“嚐鮮玩具”，但也不是“人人該買”的主流形態}。它介於小眾極客產品和大眾消費品之間，適合特定人群。

目前主流分兩類：

\begin{itemize}
  \item \textbf{橫向摺疊（書本式）}：合上是普通手機大小，展開是接近方形的小平板。代表：華為Mate X系列、三星Galaxy Z Fold系列、小米MIX Fold系列、vivo X Fold系列。主打"一部手機=手機+平板"。
  \item \textbf{縱向摺疊（翻蓋式）}：合上像粉餅盒，展開是正常尺寸手機。代表：華為Pocket系列、三星Galaxy Z Flip系列、OPPO Find N Flip、小米MIX Flip。主打"小巧便攜+外屏速覽"。
\end{itemize}

2024-2026年的關鍵進步：鉸鏈壽命普遍標稱50萬次摺疊（每天開合100次能用13年）、UTG超薄玻璃取代早期塑膠膜、摺痕肉眼可見但手感已很輕微、價格從萬元檔下探到4000-5000元檔。

\begin{tcbinfo}{一句話判斷}
如果你現在用直板機完全夠用、沒有“螢幕不夠大”的痛點，就別買摺疊屏。它是“錦上添花”的裝置，不是“雪中送炭”的裝置。

\end{tcbinfo}

\booksection{二、什麼人適合買摺疊屏？}

根據真實使用者反饋，這幾類人買摺疊屏“真香”機率高：

\begin{enumerate}
  \item \textbf{重度閱讀/文件處理者}：經常在手機上看PDF、批註文件、讀長文。展開後接近7.5英寸的內屏，看A4 PDF不用縮放，體驗質變。
  \item \textbf{經常出差/開會的商務使用者}：不想帶平板又嫌手機屏小。橫向摺疊屏展開看PPT、籤合同、雙屏對照郵件，比直板機效率高一截。
  \item \textbf{喜歡"懸停態"創作的人}：摺疊到90度放在桌上，上半屏取景下半屏按快門（自拍/延時攝影）、上半屏影片下半屏記筆記、上半屏直播下半屏看評論——這是直板機做不到的形態。
  \item \textbf{追求便攜又想要大屏的女性使用者}：縱向摺疊屏合上能塞進小包，展開是正常手機，外屏還能當化妝鏡、看通知。
\end{enumerate}

\begin{tcbwarn}{這幾類人別買}
\begin{itemize}
  \item 手機只用來刷短影片、聊微信、打電話的——摺疊屏對你毫無增益，反而更重更厚。
  \item 手小或手腕有腱鞘炎的——橫向摺疊屏展開後普遍230g+，單手握持累。
  \item 預算緊張、一部手機要用4年以上的——摺疊屏二手保值率低於直板機，維修成本高。
\end{itemize}

\end{tcbwarn}

\booksection{三、選購要點：看這6個指標}

別被“摺疊”兩個字晃花了眼，選購時重點看：

\begin{itemize}
  \item \textbf{鉸鏈型別}：水滴鉸鏈（摺痕淺、壽命長，目前主流）優於U型鉸鏈（摺痕深、早期方案）。問銷售或查評測確認是哪種。
  \item \textbf{內屏材質}：必須是UTG超薄玻璃，不要早期那種塑膠膜。UTG手感接近玻璃，塑膠膜指甲劃過會有凹陷感且易留痕。
  \item \textbf{摺痕表現}：去線下店親手展開看。2026年的合格水平是：亮屏時正面幾乎看不出，斜光下能看出一條淺溝，手指劃過有輕微凹陷但不割手。
  \item \textbf{外屏尺寸}：橫向摺疊屏的外屏越大越好（最好6.4英寸以上），合上能當正常手機用；縱向摺疊屏的外屏越大越實用（4英寸以上能回訊息、看導航）。
  \item \textbf{重量與厚度}：橫向摺疊展開後230-260g是常態，低於230g算優秀；合上厚度12-14mm，低於12mm算薄。縱向摺疊合上後能輕鬆放進襯衫口袋才算合格。
  \item \textbf{售後與保修}：重點看內屏碎裂的維修價格。有些品牌首年免費換一次內屏，有些自費要2000-4000元。這是摺疊屏最大的隱性成本。
\end{itemize}

\begin{tcbpractice}{線下體驗清單}
\begin{enumerate}
  \item 親手開合10次，感受鉸鏈阻尼是否順滑、有沒有異響。
  \item 展開後斜對著燈光看摺痕，再正面亮屏看。
  \item 單手握持展開狀態30秒，感受重量是否能接受。
  \item 試用懸停態：折90度放桌上，試試拍照和看影片。
  \item 問清楚內屏單獨維修價格和保修政策。
\end{enumerate}

\end{tcbpractice}

\booksection{四、摺疊屏獨有的使用技巧}

買回來之後，這幾個玩法是直板機體驗不到的：

\booksubsection{分屏與多工}

橫向摺疊屏展開後，大部分系統支援上下/左右分屏。典型場景：上半屏看網課影片，下半屏記筆記；左邊微信右邊文件；上面地圖導航下面音樂。比直板機分屏實用得多，因為螢幕夠大，每半邊都不擠。

\booksubsection{懸停態（90度摺疊）}

這是摺疊屏的“殺手鐧”用法：

\begin{itemize}
  \item \textbf{自拍/合影}：放桌上折90度，上半屏取景，下半屏按快門，不用找人幫忙拍。
  \item \textbf{延時攝影/夜景}：同樣放桌上穩定拍攝，省三腳架。
  \item \textbf{視訊通話}：解放雙手，對方看你角度也更自然。
  \item \textbf{看影片}：放桌上當小電視，飛機/高鐵小桌板上特別實用。
  \item \textbf{外屏預覽+內屏拍攝}：用後置攝像頭拍自己時，外屏能實時顯示畫面，構圖更準。
\end{itemize}

\booksubsection{外屏速覽}

縱向摺疊屏合上後，外屏可以：看時間/通知/天氣、回微信（語音轉文字或快捷回覆）、刷公交卡、看導航、當化妝鏡（呼叫前置攝像頭）。不用每次都展開，省電又省事。

\booksubsection{電腦模式}

部分橫向摺疊屏（華為、三星）展開後可切換“電腦模式”，外接顯示器鍵鼠後像臺簡易PC，臨時處理文件、回郵件夠用。不是主力生產力工具，但應急很香。

\booksection{五、保養要點：讓摺疊屏活得久一點}

摺疊屏最嬌貴的就是內屏和鉸鏈，這幾條習慣能顯著延長壽命：

\begin{itemize}
  \item \textbf{別在低溫下強行開合}：0℃以下UTG玻璃變脆，開合前先揣懷裡暖幾分鐘。這是官方說明書都會寫但沒人看的重點。
  \item \textbf{避免指甲/硬物接觸內屏}：內屏雖有保護層但比外屏軟，用指腹操作，別用指甲戳。
  \item \textbf{合上時確認屏內無異物}：一粒沙子夾在中間合上，可能直接壓壞內屏。合上前掃一眼。
  \item \textbf{不要貼普通鋼化膜}：內屏只能用官方專用膜或UTG保護膜，普通鋼化膜會影響摺疊和觸控。建議去官方售後貼膜。
  \item \textbf{鉸鏈別進灰塵液體}：雖然多數標稱IPX8防水，但防塵等級普遍較低。沙灘、工地等環境要小心。
  \item \textbf{開合力度均勻}：別單手暴力甩開，雙手從中間勻速展開最穩。
\end{itemize}

\textbf{內屏損壞的高發場景}\par

\begin{itemize}
  \item 合上時夾住耳機線/充電線——直接壓裂內屏。
  \item 放包裡和鑰匙混裝——外屏沒事，內屏被鑰匙硌出裂痕。
  \item 冬天從冷處進暖房立刻開合——溫差導致玻璃應力變化。
  \item 用了一年後覺得"摺痕變深了"去撕原廠膜——撕膜會連內屏一起損壞。
\end{itemize}

\booksection{六、摺疊屏的“坑”和心理預期}

買之前心裡要有這幾本賬：

\begin{itemize}
  \item \textbf{摺痕會一直在}：2026年的技術做不到"完全無痕"，只能做到"亮屏正面幾乎看不出"。如果你是完美主義者，每天盯著摺痕難受，別買。
  \item \textbf{更重更厚}：橫向摺疊屏展開後普遍230g+，單手握持比直板機累；合上後厚度12-14mm，褲兜裡更鼓。
  \item \textbf{續航普遍一般}：因為要塞鉸鏈機構，電池空間被壓縮，同價位續航通常不如直板機。
  \item \textbf{保值率低}：一年後二手價通常是新機的50-60\%，比直板機跌得快。
  \item \textbf{維修貴}：內屏總成維修價2000-4000元，鉸鏈維修500-1500元。買之前查清楚保修政策。
  \item \textbf{App適配參差}：主流App已適配大屏，但小眾App可能只是"拉伸放大"，體驗一般。
\end{itemize}

摺疊屏不是“更好的手機”，而是“另一種形態的手機”。買它是因為你需要它獨特的形態，不是因為它“更先進”。

—— 關於摺疊屏的選購哲學

\booksection{七、主流品牌怎麼選（2026年）}

不寫具體型號（會過時），只寫品牌調性，幫你快速定位：

\begin{itemize}
  \item \textbf{華為Mate X系列}：橫向摺疊標杆，鴻蒙生態協同強，外屏比例接近正常手機，合上能當主力機用。價格高但保值率相對最好。
  \item \textbf{三星Galaxy Z Fold/Flip}：摺疊屏鼻祖，鉸鏈調校成熟，Z Flip是縱向摺疊的顏值擔當。國內售後網點較少是短板。
  \item \textbf{小米MIX Fold/Flip}：價效比路線，配置堆料足，價格相對親民。適合想嚐鮮又不想花太多的使用者。
  \item \textbf{vivo X Fold系列}：橫向摺疊裡外屏比例做得最接近正常手機的，合上當直板機用體驗好。
  \item \textbf{OPPO Find N系列}：橫向摺疊裡尺寸偏緊湊的，單手可握持，適合覺得大摺疊太重的使用者。
  \item \textbf{華為Pocket系列}：縱向摺疊裡外屏最大最實用的，化妝鏡、防曬檢測等女性向功能做得細緻。
\end{itemize}

選購建議：預算充足追求生態選華為Mate X；想要最成熟鉸鏈選三星Z Fold；價效比首選小米/vivo；縱向摺疊看華為Pocket和三星Z Flip。

\begin{tcbinfo}{最後提醒}
摺疊屏技術每年都在迭代，本附錄寫於2026年。買之前務必查最新評測和使用者反饋，重點關注：鉸鏈壽命實測、摺痕表現、內屏維修政策這三項。線下親手摸永遠比看評測重要。

\end{tcbinfo}
